\section{Introduction and Motivation}
\label{sec:introduction}

Emotion detection asks which emotions a piece of text expresses. Dialogue makes this harder, as a movie subtitle is short, carries almost no context, and often conveys emotion implicitly through irony, a question, or a reaction to something the reader cannot see. A line can also express several emotions at once; ``You did what?'' may be read as both surprise and anger. The task is therefore multi-label, in that a classifier must decide, for each of a fixed set of emotions, whether it is present, so partially correct answers are common and rare emotions are hard both to learn and to evaluate.

Large language models (LLMs) offer two ways to approach this task. They can be prompted to label emotions directly, from a description of the task and optionally a few solved examples, or, when labelled data is available, adapted to it by training a small number of additional weights. Whether either works for a fine-grained, multi-label task beyond English is less clear: LLMs see far less text in most other languages during pretraining, and labelled data outside English is scarce. In the XED dataset \parencite{ohman2020xed}, only English and Finnish were annotated by people, and the labels for other languages were projected onto aligned subtitles, so any evaluation in these languages inherits their noise.

In this project, we study how well open-weight LLMs classify emotions in English and in Croatian, our assigned language, comparing prompt engineering, with zero or five examples in the prompt, and parameter-efficient fine-tuning (PEFT). Training on English alone and evaluating on Croatian also shows how much of what is learned transfers across languages. Croatian falls in band~3 of the project's five resource bands (``mid-low'', roughly 4,400--6,000 annotated lines; Table~\ref{tab:bands}). No member of our group reads Croatian, so Croatian outputs are judged only against XED's projected labels. We:

\begin{itemize}
    \item build a fixed, leakage-controlled split of matched English--Croatian pairs;
    \item fine-tune our primary model, EuroLLM-9B, on English, Croatian and both, with one configuration selected on the development set, and retrain on smaller Croatian samples to see how much target-language data fine-tuning needs;
    \item prompt EuroLLM-9B and a second model, Qwen3.5-9B, zero-shot and five-shot, and measure sensitivity to prompt wording; and
    \item evaluate every condition with multi-label metrics and bootstrap confidence intervals, and analyse concrete errors, separating model errors from those that may stem from subtitle alignment or label projection.
\end{itemize}

Section~\ref{sec:background} introduces multi-label emotion classification, XED and its projected labels, and the two adaptation approaches. Section~\ref{sec:data} describes the data and split, Section~\ref{sec:methodology} the models, fine-tuning and prompts, and Section~\ref{sec:evaluation} the evaluation protocol. Sections~\ref{sec:results} and~\ref{sec:analysis} present and discuss the results, and Section~\ref{sec:conclusion} concludes. The code and configuration files are available at \url{https://github.com/squidistaken/large-language-models}.
